\chapter{Results from the literature}\label{chapter-literature}
Alpöge's proof uses about one hundred and fifty published results, and each
of them is proved in Alpöge's archive. Many are proved in the form of the
published theorem; several are special cases or weaker statements that
suffice where they are used, and one entry of the list is the true
proposition. This chapter states in full the results that the other chapters
quote, and the ones whose proved form differs from the published theorem in a
way a reader should know about. The remaining results are listed, with their
published sources, in Theorems~\ref{lit-bundle-unitary}
and~\ref{lit-bundle-theta}.

\section{How the results enter the proof}\label{section-lit-inputs}

\begin{remark}\label{lit-open-inputs}
The results of this chapter are proved under three hypotheses: two copies of
the statement of Definition~\ref{gal-galois-representations}, and the statement
of Theorem~\ref{gal-llc-exists}. The two copies carry the same statement under
two names, one of which refers to Caraiani's local-global compatibility.
Only the results concerning the local Langlands correspondence, the Galois
representations attached to polarized cuspidal representations, local-global
compatibility, and the $\ell$-adic characters attached to Hecke characters
above $\ell$ use these hypotheses; all others are proved without them.
Chapter~\ref{chapter-galois-representations} proves both hypotheses, in
Theorems~\ref{gal-galois-representations-any-rc} and~\ref{gal-llc-exists},
and Chapter~\ref{chapter-assembly} supplies them, so every result of this
chapter holds unconditionally.
\end{remark}
\leanref{Development/AutomorphicArithmetic/ThetaCorrespondence/ThetaPairLifts/ThirdPlaceSeries/EveryCharacterOfLiteratureOfVWin.lean}{UpperTheory.LiteratureTheorems.OpenLiteratureStatements}

\begin{theorem}\label{lit-bundle-unitary}
Assume the property of Definition~\ref{gal-galois-representations}
and the statement of Theorem~\ref{gal-llc-exists}. Then the following results
hold. Those stated in this book are cited by their number; for the others the
published source is cited.
\begin{enumerate}
\item Real unitary groups $U(p,q)$: Theorem~\ref{lit-casselman-wallach};
  Harish-Chandra's finiteness theorem for automorphic forms on unitary groups
  invariant under a given open subgroup of the finite adelic points, with
  given $K_\infty$-types and infinitesimal character \cite{HC68}; the
  description of the irreducible unitary representations of $U(p,q)$ by their
  Harish-Chandra modules \cite{HC53}; analyticity of $K$-finite matrix
  coefficients and density of the $(\mathfrak g,K)$-stable subspaces of
  $K$-finite vectors \cite{HC53}; the existence of the discrete series of
  $U(p,q)$, their classification by Harish-Chandra parameters, and their
  covering by representations with prescribed lowest $K$-type, Casimir
  eigenvalue and highest weights \cite{HC65,HC66,Sch75,HS75}; the formal degree of
  a discrete series representation in terms of the compact dual
  \cite{HC66,HII08}; and Schur orthogonality for square-integrable
  representations of unimodular groups \cite[Chapter~14]{Dix77}.
\item General linear groups over a nonarchimedean local field: the finiteness
  theorems of Bernstein and Zelevinsky (irreducible smooth representations are
  admissible, Jacquet modules and parabolic induction preserve finite length,
  induction commutes with contragredients) \cite{BZ76,BZ77}; Zelevinsky's
  isomorphism of the Grothendieck ring with a polynomial ring on segments
  \cite[\S7]{Zel80}; Bernstein's theorem that a product of irreducible
  unitarisable representations is irreducible \cite{Ber84}; the filtration of
  the geometric lemma and the semisimplification of its pieces
  \cite[2.12]{BZ77}; irreducibility of a product of two irreducible
  representations whose cuspidal supports are not related by unramified
  twists, or not by twists by $|\det|^j$ with $|j|\le1$ \cite{BZ77,Zel80};
  injectivity, and bijectivity, of the map from depth-zero invariants to the
  invariants of the Jacquet module, and the fact that depth zero passes to and
  from the cuspidal support \cite{MP94,MP96};
  Tadić's description of the generic unitary representations \cite{Tad86};
  existence of the cosocle of a standard module \cite{Sil78,Zel80}; and Green's
  description of the cuspidal characters of $GL_d(\mathbf F_g)$
  (Theorem~\ref{lit-green-cuspidal} and
  Corollary~\ref{lit-green-labelled}).
\item Theta integrals and Siegel Eisenstein series:
  Theorems~\ref{lit-weil-siegel}, \ref{lit-siegel-weil}
  and~\ref{lit-doubling}; Siegel's formula expressing the mean of
  representation numbers of a Hermitian form by a product of local densities
  in the range of absolute convergence, with the finiteness of the volume of
  the automorphic quotient \cite{Wei65,BHC62}; the Siegel-Weil formula of
  Theorem~\ref{lit-siegel-weil} in the Schrödinger model \cite{Ich07}; the
  meromorphic continuation of Siegel Eisenstein series, which is part~(1) of
  Theorem~\ref{lit-doubling} \cite{Lan76,MW95}; the unramified computation of
  the doubling integral for $GL_n$ \cite{GPSR87,LR05}; the existence of the
  Weil representation of a unitary dual pair \cite{Kud94,HKS96}; and the entry
  described in Remark~\ref{lit-vacuous-slot}.
\item Automorphic forms: rapid decay of cusp forms on unitary groups
  \cite{HC68,MW95}; the trace class property of the cuspidal spectrum
  \cite{GGPS69}; the factorisation theorem of Dixmier and Malliavin
  \cite{DM78}; Flath's factorisation of irreducible unitary representations of
  adelic unitary groups \cite{Fla79}; strong approximation for special unitary
  groups \cite{Kne65,PR94}; finite dimensionality of spaces of automorphic forms
  on $GL_a$ over a totally complex field with given invariance under an open
  subgroup, $K_\infty$-type and infinitesimal character \cite{HC68,BJ79};
  Theorems~\ref{lit-jacquet-shalika} and~\ref{lit-clozel-cohomological};
  Clozel's purity lemma for regular algebraic cuspidal representations
  \cite[Lemme~4.9]{Clo90}; and base change from definite unitary groups
  \cite{Lab11}, which is deduced from Theorem~\ref{gal-base-change-definite}.
\item Galois representations and class field theory: global Artin reciprocity
  and the existence of ray class fields and of the Artin characters attached
  to ray class characters \cite[Chapter~VI]{Neu99}; the $\ell$-adic characters
  attached to algebraic Hecke characters, with their Hodge-Tate weights and
  Weil-Deligne parameters at places above $\ell$ \cite{Ser68,Fon94}; the
  Galois representations attached to regular algebraic polarized cuspidal
  representations and their local-global compatibility above $\ell$
  \cite{CH13,BLGGT14,Car14}, deduced from the property of
  Definition~\ref{gal-galois-representations} and Theorem~\ref{gal-llc-exists};
  Berger's theorem that de Rham representations are potentially semistable
  \cite{Ber02}; Theorem~\ref{lit-bhatt-scholze};
  Theorem~\ref{lit-chebotarev-degree-one} and a second form of the
  Chebotarev density theorem \cite{Che26,Neu99}; the existence, for each prime
  $p\equiv3\pmod 4$ with $p>3$, of an elliptic curve with complex
  multiplication by the integers of $\mathbf Q(\sqrt{-p})$, good reduction
  outside $p$ and minimal discriminant $-p^3$ \cite{Gro80}; Weber's
  description of ray class fields of imaginary quadratic fields
  \cite[Theorem~II.5.6]{Sil94}; and local Tate duality, the finiteness of
  local Galois cohomology and the local Euler characteristic formula
  \cite[Chapter~VII]{NSW08}.
\item Three further collections. The first contains: the local results of
  item~(2), together with the characterization of essentially
  square-integrable and of tempered representations by segments, the
  Langlands classification by multisegments and Zelevinsky's description of
  the generic representations \cite[\S9]{Zel80}; the local converse theorem
  for cuspidal representations of $GL_n$ \cite{Hen93}; Brauer's theorem on
  induced characters in degree zero \cite{Bra47}; Chevalley's theorem on
  congruence subgroups of units \cite{Che51}; the statement of
  Theorem~\ref{gal-llc-exists}; the global results of items~(4) and~(5),
  together with the discreteness of the cuspidal spectrum of $GL_a$ and the
  boundedness of its cusp forms \cite{GGPS69,MW95} and the description of the
  contragredient of a cuspidal representation by the transpose inverse
  \cite{GK75,JS81b}; the statement of
  Definition~\ref{gal-galois-representations}; Poitou-Tate duality and the
  global finiteness and Euler characteristic formulas
  \cite[Chapter~VIII]{NSW08}; the orthogonality of unramified classes under
  local duality \cite[Chapter~VII]{NSW08}; the formula of Greenberg and Wiles
  \cite[Proposition~1.6]{Wil95}; the local results of Bloch and Kato
  \cite[\S3]{BK90}; five statements of Alpöge's on lattices with complex
  multiplication and on denominators and reflection identities of zeta and
  $L$-values, related to \cite{Shi71,Dam70,DR80}; and the functional equation
  of Hecke $L$-functions of CM type \cite{Hec20,Tat67}. The second contains
  the Weil representation, the existence and classification of discrete
  series, Harish-Chandra analyticity, the trace class property, the
  factorisations of Dixmier-Malliavin and Flath,
  Theorems~\ref{lit-weil-siegel} and~\ref{lit-doubling}, the unramified
  doubling integral, the continuation of partial Hecke $L$-functions
  \cite{Hec20,Tat67}, and the fact that a de Rham representation of dimension
  $a$ has $a$ Hodge-Tate weights counted with multiplicity, for the
  representations attached to the data of Definition~\ref{th-definite-datum}
  \cite{Fon94}. The third contains Fontaine's theory of the ring
  $B_{\mathrm{cris}}$ \cite{Fon94a}, global Artin reciprocity, the local
  results of Bloch and Kato, local Tate duality, finiteness and Euler
  characteristic over finite extensions of $\mathbf Q_\ell$, and the integral
  form of the theorem of Bhatt and Scholze \cite{BS23}.
\end{enumerate}
\end{theorem}
\leanref{Development/AutomorphicArithmetic/ThetaCorrespondence/ThetaPairLifts/ThirdPlaceSeries/EveryCharacterOfLiteratureOfVWin.lean}{UpperTheory.LiteratureTheorems.published_results_eigen}
\leanref{Development/AutomorphicArithmetic/ThetaCorrespondence/ThetaPairLifts/ThirdPlaceSeries/BernsteinUnitaryProductIrreducible.lean}{GlobalTheory.Fc.S0.UnitaryGroupThirdPlaceHypotheses}

\begin{theorem}\label{lit-bundle-theta}
The following results hold. Those stated in this book are cited by their
number; for the others the published source is cited.
\begin{enumerate}
\item Finite and $p$-adic groups: for each monic squarefree polynomial $f$ of
  degree $N$ over $\mathbf F_g$ with nonzero constant term, an irreducible
  generic character of $GL_N(\mathbf F_g)$ of degree
  $\prod_{i=1}^N(g^i-1)/|T_f|$, where $T_f$ is the torus attached to $f$,
  distinct for distinct $f$ \cite{Gre55}; the existence of cuspidal
  representations of $GL_n$ over a finite field and the description of their
  restriction to the mirabolic subgroup \cite{Gel70,Zel81}; Bernstein's
  irreducibility theorem, Zelevinsky's polynomial ring and the finiteness
  theorems of Bernstein and Zelevinsky, as in item~(2) of
  Theorem~\ref{lit-bundle-unitary}; the theorem of Borel and Casselman that
  the Iwahori invariants of an admissible representation of $GL_n$ map
  bijectively to the invariants of its Jacquet module \cite{Bor76}; the local
  classification of Hermitian forms over a quadratic extension of local
  fields by their determinant \cite{Jac62}; Kudla's splitting of the Weil
  representation for split unitary groups \cite{Kud94}; and the Stone-von
  Neumann theorem \cite[Chapitre~2]{MVW87}.
\item Real groups: Theorems~\ref{lit-casselman-wallach}
  and~\ref{lit-howe-duality-real}; Weyl's branching law
  \cite[Chapter~8]{GW09}; Howe duality for
  $GL_n(\mathbf C)\times GL_{n'}(\mathbf C)$ and $U(n)\times U(n')$ acting on
  polynomials \cite{How89b}; two forms of the real metaplectic representation
  \cite{Wei64}; $K$-finite vectors of irreducible unitary representations of
  $U(p,q)$ and the globalisation of $(\mathfrak g,K)$-modules \cite{HC53}; the
  existence of the discrete series of $U(p,q)$, their lowest $K$-types,
  uniqueness for a given Harish-Chandra parameter, and the lowest weight
  description of the holomorphic discrete series
  \cite{HC56,HC66,Sch75}; their formal degrees \cite{HC66,HII08};
  Harish-Chandra modules \cite{HC53}; and Schur orthogonality
  \cite[Chapter~14]{Dix77}.
\item Global theory: Theorem~\ref{lit-linnik}; the main theorem of complex
  multiplication for $j$-invariants of imaginary quadratic fields and the
  transitivity of the Galois action on them \cite[Chapter~II]{Sil94}; quadratic
  class field theory \cite[Chapter~VI]{Neu99}; the Hasse principle for
  Hermitian forms and its local-global form \cite{Lan36,Jac62}; boundedness,
  decay estimates and rapid decay of cusp forms \cite{HC68,MW95}; trace class
  of the cuspidal spectrum \cite{GGPS69}; the factorisations of
  Dixmier-Malliavin and Flath \cite{DM78,Fla79}; strong approximation for
  special unitary groups \cite{Kne65,PR94}; Siegel sets for unitary groups
  \cite{Bor69,BHC62}; Harish-Chandra's finiteness theorem \cite{HC68};
  holomorphy of Hecke $L$-functions of unitary Hecke characters other than
  powers of the norm, with nonvanishing on $\operatorname{Re}s\ge1$, also for
  partial $L$-functions \cite{Hec20,Tat67,Lan94}; the bound
  $\bigl|\frac{L'}{L}(1,\chi)-\sum_\rho\frac1{1-\rho}\bigr|\le C\log(3q)$ for
  primitive Dirichlet characters $\chi$ of conductor $q>1$, the sum running
  over the zeros $\rho$ of $L(s,\chi)$ with $0<\operatorname{Re}\rho<1$ and
  $|\operatorname{Im}\rho|\le1$, counted with multiplicity
  \cite[Chapter~16]{Dav00}; and Radford's theorem that the shuffle algebra over
  a field of characteristic zero is the polynomial algebra on Lyndon words
  \cite{Rad79}.
\item Theta correspondence: Theorems~\ref{lit-weil-siegel}
  and~\ref{lit-doubling}; the Rallis inner product formula in the convergent
  range \cite{Ral84,Li92} (compare Theorem~\ref{th-rallis-inner-product}); the
  unramified doubling integral for $GL_n$ \cite{GPSR87,LR05}; the invariance of
  theta functions under rational points, in an explicit form \cite{Wei64}; and
  two forms of the existence of the Weil representation of a unitary dual
  pair \cite{Kud94,HKS96}.
\end{enumerate}
\end{theorem}
\leanref{Development/AutomorphicArithmetic/ThetaCorrespondence/ThetaPairLifts/ThirdPlaceSeries/EveryCharacterOfLiteratureOfVWin.lean}{UpperTheory.LiteratureTheorems.published_results_first}
\leanref{Development/AutomorphicArithmetic/ThetaCorrespondence/ThetaPairLifts/LocalTypeDictionary/BernsteinUnitaryProductIrreducible.lean}{GlobalTheory.FunPr.ThetaLocalInputs}

\section{Representations of real unitary groups}\label{section-lit-real}

\begin{theorem}[Casselman, Wallach \cite{Cas89,Wal92}]\label{lit-casselman-wallach}
Let $n\ge1$ and let $J$ be a diagonal $n\times n$ matrix with entries $\pm1$.
Let $G=\{g\in GL_n(\mathbf C):g^*Jg=J\}$, let $K=G\cap U(n)$ and let
$\mathfrak g=\{X\in M_n(\mathbf C):X^*J+JX=0\}$. For $g\in M_n(\mathbf C)$
put $\|g\|^2=\sum_{a,b}|g_{ab}|^2$.

Let $(\sigma,H)$ be an irreducible unitary representation of $G$ on a Hilbert
space. Let $H_K\subset H$ be its space of $K$-finite vectors and
$H^\infty\subset H$ its space of smooth vectors, and for a finite set $S$ of
words $X_1\cdots X_k$ in elements of $\mathfrak g$ put
$\|w\|_S=\sum_{X_1\cdots X_k\in S}\|d\sigma(X_1)\cdots d\sigma(X_k)w\|$.

Let $E$ be a complete Hausdorff locally convex space whose topology is given
by a sequence of seminorms $(p_i)_{i\ge0}$. Let $\pi$ be a representation of
$G$ on $E$ by continuous operators such that $g\mapsto\pi(g)v$ is continuous
for every $v$, and let $d\pi(X)$, $X\in\mathfrak g$, be continuous operators
with $\frac1t(\pi(\exp tX)v-v)\to d\pi(X)v$ as $t\to0$ for every $v\in E$.
Assume that $\pi$ has moderate growth: for every $i$ there are a continuous
seminorm $q$ on $E$ and $N\ge0$ with
$p_i(\pi(g)v)\le(1+\|g\|^2)^Nq(v)$ for all $g\in G$ and $v\in E$.

Let $T_0:H_K\to E$ be a map of $(\mathfrak g,K)$-modules. Then:
\begin{enumerate}
\item there is a linear map $T:H^\infty\to E$ that extends $T_0$, commutes
  with the action of $G$, and is continuous in the sense that for every $i$
  there are a finite set $S$ and a constant $C$ with
  $p_i(Tw)\le C\|w\|_S$ for all $w\in H^\infty$;
\item any two maps with the properties of~(1) are equal;
\item $T(H^\infty)$ is closed in $E$;
\item if $T_0\neq0$, then for every finite set $S$ of words there is a
  continuous seminorm $q$ on $E$ with $\|w\|_S\le q(Tw)$ for all
  $w\in H^\infty$.
\end{enumerate}
\end{theorem}
\leanref{Development/AbelianVarietyTheory/PublishedTheorems/AutomorphicForms₂/HasDRepIffHasDerivAt.lean}{GlobalTheory.Published.CasselmanWallach}
\leanref{Development/WeilBound/CoreStatements/PublishedTheorems/CasselmanWallach/FinalB.lean}{PositivityTheory.CWB.casselmanWallach}

\begin{theorem}[Howe \cite{How89}]\label{lit-howe-duality-real}
Let $E/F$ be a CM extension of number fields with complex conjugation $c$, and
let $\eta$ be the quadratic character of $\mathbf A_F^\times/F^\times$
attached to it. Let $X$ be a nondegenerate Hermitian space of dimension $n$
over $E$ and $W$ a nondegenerate skew-Hermitian space of dimension $d$ over
$E$. Let $\psi$ be a standard additive character of $\mathbf A_F/F$. Let
$\chi_X,\chi_W$ be Hecke characters of $\mathbf A_E^\times/E^\times$ with
$\chi_X|_{\mathbf A_F^\times}=\eta^n$ and
$\chi_W|_{\mathbf A_F^\times}=\eta^d$, and let $\omega$ be a Weil
representation of $U(X)(\mathbf A_F)\times U(W)(\mathbf A_F)$ attached to
$(\psi,\chi_X,\chi_W)$, realised on Schwartz functions on
$\mathbf A_F^{n\times d}$.

For each archimedean place $w$ of $F$ choose an embedding $\tau_w$ of $E$
over $w$ and Cartan involutions $\theta^X_w$ of $X\otimes_{\tau_w}\mathbf C$
and $\theta^W_w$ of $W\otimes_{\tau_w}\mathbf C$, the second for the
Hermitian form $i$ times that of $W$. Let $K_X\subset U(X)(F_\infty)$ and
$K_W\subset U(W)(F_\infty)$ be the subgroups commuting with these
involutions, and let $\mathfrak u(X)$ and $\mathfrak u(W)$ be the Lie algebras
of $U(X)(F_\infty)$ and $U(W)(F_\infty)$. Let $\mathcal P$ be the Fock space:
the smallest space of Schwartz functions on $F_\infty^{n\times d}$ that
contains the Gaussian vectors attached to the involutions and is stable under
the derivatives of the archimedean Heisenberg action.

Call a subspace $N\subset\mathcal P$ a \emph{theta subspace} if it is stable under
$d\omega(\mathfrak u(W))$ and $\omega(K_W)$ and contains $d\omega(Y)f$ and
$\omega(k)f-f$ for all $f\in\mathcal P$, $Y\in\mathfrak u(X)$ and
$k\in K_X$. Then:
\begin{enumerate}
\item there is $L$ such that every strictly increasing chain of theta
  subspaces has at most $L$ members;
\item if $N,N'$ are theta subspaces different from $\mathcal P$, then
  $N+N'\neq\mathcal P$.
\end{enumerate}
Equivalently, the largest quotient of $\mathcal P$ on which
$(\mathfrak u(X),K_X)$ acts trivially is a $(\mathfrak u(W),K_W)$-module of
finite length with at most one irreducible quotient.
\end{theorem}
\leanref{Development/AbelianVarietyTheory/PublishedTheorems/Divisors₂/HoweDualityReal.lean}{GlobalTheory.Published.HoweDualityReal}
\leanref{Development/AbelianVarietyTheory/PublishedTheorems/Divisors₂/HoweDualityReal.lean}{GlobalTheory.Published.ThetaSubInf}
\leanref{Development/AbelianVarietyTheory/Proofs/DivisorProofs/HoweDualityRealHolds.lean}{GlobalTheory.howeDualityReal_holds}

\section{Cuspidal characters of finite general linear groups}\label{section-lit-green}

\begin{theorem}[Green \cite{Gre55}]\label{lit-green-cuspidal}
Let $g$ be a prime, let $d\ge1$, and let $K$ be a field extension of
$\mathbf F_g$ of degree $d$. Call a character $\theta$ of $K^\times$ regular if
the characters $\theta^{g^i}$, $0\le i<d$, are distinct. Call an irreducible
character $\chi$ of $GL_d(\mathbf F_g)$ cuspidal if, for every
$0<a<d$, the unipotent radical of the standard parabolic subgroup of type
$(a,d-a)$ has no nonzero fixed vector in the representation with
character $\chi$. Then there is a map $\theta\mapsto\tau_\theta$ from
characters of $K^\times$ to class functions on $GL_d(\mathbf F_g)$ such that:
\begin{enumerate}
\item for regular $\theta$, $\tau_\theta$ is a cuspidal irreducible character
  of degree $\prod_{i=1}^{d-1}(g^i-1)$, and
  \[
  \tau_\theta(y)=(-1)^{d-1}\sum_{i=0}^{d-1}\theta(z^{g^i})
  \]
  for every $y\in GL_d(\mathbf F_g)$ whose characteristic polynomial is
  irreducible and every root $z\in K^\times$ of that polynomial;
\item for regular $\theta,\theta'$, $\tau_\theta=\tau_{\theta'}$ if and only if
  $\theta'=\theta^{g^i}$ for some $0\le i<d$;
\item every cuspidal irreducible character of $GL_d(\mathbf F_g)$ is
  $\tau_\theta$ for some regular $\theta$.
\end{enumerate}
\end{theorem}
\leanref{Development/AbelianVarietyTheory/PublishedTheorems/GaloisCohomology₂/GreenRegularCharacters.lean}{GlobalTheory.Published.GreenRegularCharacters}
\leanref{Development/WeilBound/CoreStatements/StandardPublishedResults/StandardTheoremStatements/GreenRegularCharacters.lean}{PositivityTheory.Stage1.PM.greenRegularCharacters}

\begin{definition}\label{lit-labelling}
Let $g$ be a prime. A \emph{labelling} of $\bar{\mathbf F}_g$ consists of an
injective homomorphism $t:\bar{\mathbf F}_g^\times\to\mathbf C^\times$ and, for
each $d\ge1$, a generator $\gamma_d$ of $\mathbf F_{g^d}^\times$, such that
$\gamma_{de}^{(g^{de}-1)/(g^d-1)}=\gamma_d$ for all $d,e\ge1$. For
$x\in\bar{\mathbf F}_g^\times$ and $d\ge1$, let $\theta_{x,d}$ be the
character of $\mathbf F_{g^d}^\times$ with $\theta_{x,d}(\gamma_d^m)=t(x)^m$.
\end{definition}
\leanref{Development/AbelianVarietyTheory/Definitions/NonGaloisDescent/LabelDatum.lean}{GlobalTheory.GLq.LabelDatum}
\leanref{Development/AbelianVarietyTheory/Definitions/NonGaloisDescent/LabelDatum.lean}{GlobalTheory.GLq.thetaOf}

\begin{corollary}\label{lit-green-labelled}
Let $g$ be a prime, let $d\ge1$, and fix a labelling of $\bar{\mathbf F}_g$
(Definition~\ref{lit-labelling}). Let $\mathcal P_d$ be the set of monic
irreducible polynomials of degree $d$ over $\mathbf F_g$ with nonzero constant
term. There is a map $P\mapsto\tau_P$ from $\mathcal P_d$ to class functions
on $GL_d(\mathbf F_g)$ such that:
\begin{enumerate}
\item each $\tau_P$ is a cuspidal irreducible character of degree
  $\prod_{i=1}^{d-1}(g^i-1)$, and
  $\tau_P(y)=(-1)^{d-1}\sum_{i=0}^{d-1}\theta_{x,d}(z^{g^i})$ for every root
  $x$ of $P$, every $y\in GL_d(\mathbf F_g)$ with irreducible characteristic
  polynomial and every root $z$ of that polynomial;
\item $P\mapsto\tau_P$ is injective;
\item every cuspidal irreducible character of $GL_d(\mathbf F_g)$ is $\tau_P$
  for some $P\in\mathcal P_d$.
\end{enumerate}
\end{corollary}
\begin{proof}
Apply Theorem~\ref{lit-green-cuspidal} with $K=\mathbf F_{g^d}$. For
$P\in\mathcal P_d$ choose a root $x_P$ and put
$\tau_P=\tau_{\theta_{x_P,d}}$. The roots of $P$ are
$x_P^{g^i}$, $0\le i<d$, and they are distinct. Since $t$ is injective,
$\theta_{x,d}^{g^i}=\theta_{x^{g^i},d}$ and $\theta_{x,d}=\theta_{x',d}$ only
when $x=x'$. So $\theta_{x_P,d}$ is regular, and
Theorem~\ref{lit-green-cuspidal}(1) gives the degree and cuspidality.

For the character formula, let $y$ have irreducible characteristic
polynomial $f$ and let $z$ be a root of $f$ in $\bar{\mathbf F}_g$. Then $f$
has degree $d$ and nonzero constant term, so $z\in\mathbf F_{g^d}^\times$ and
the formula of Theorem~\ref{lit-green-cuspidal}(1) applies. Replacing $x_P$ by
another root $x_P^{g^j}$ replaces $\theta$ by $\theta^{g^j}$, which permutes
the terms $\theta(z^{g^i})$ cyclically, so the sum does not depend on the root.

If $\tau_P=\tau_{P'}$, then by Theorem~\ref{lit-green-cuspidal}(2)
$\theta_{x_{P'},d}=\theta_{x_P,d}^{g^i}=\theta_{x_P^{g^i},d}$ for some $i$, so
$x_{P'}=x_P^{g^i}$ is a root of $P$ and $P'=P$. Finally, a cuspidal character
is $\tau_\theta$ with $\theta$ regular by
Theorem~\ref{lit-green-cuspidal}(3). Every character of the cyclic group
$\mathbf F_{g^d}^\times$ is $\theta_{x,d}$ with $x=t^{-1}(\theta(\gamma_d))$,
which exists because $t$ restricts to an isomorphism from the
$(g^d-1)$-st roots of unity in $\bar{\mathbf F}_g$ onto those in $\mathbf C$.
Regularity of $\theta$ says that $x,x^g,\ldots,x^{g^{d-1}}$ are distinct, so
$x$ has degree $d$ over $\mathbf F_g$ and its minimal polynomial $P$ lies in
$\mathcal P_d$, with $\tau_P=\tau_\theta$ by~(2).
\end{proof}
\leanref{Development/AbelianVarietyTheory/Proofs/InertiaEigenvalues/SumRangeEqRoots.lean}{GlobalTheory.GreenBridge.green_of_regular}
\leanref{Development/AbelianVarietyTheory/Definitions/NonGaloisDescent/LabelDatum.lean}{GlobalTheory.GLq.Green}

\section{Theta integrals and Siegel Eisenstein series}\label{section-lit-theta}

\begin{definition}\label{lit-weil-datum}
Let $\mathcal D$ be a definite-side datum
(Definition~\ref{th-definite-datum}), with CM field $E$, totally real
subfield $F$ and quadratic character $\eta$ of $E/F$. Let $S$ and $T$ be
nondegenerate Hermitian spaces over $E$ of dimensions $s$ and $t$, and let
$\mathbb W_{S,T}$ be the space of $t\times s$ matrices over
$\mathbf A_E=\mathbf A_F\otimes_FE$, with the symplectic $\mathbf A_F$-form
built from the two Hermitian forms and the Heisenberg action
$(w_0\cdot f)(w)=\psi(\tfrac12\langle w,w_0\rangle)f(w+w_0)$.
\begin{enumerate}
\item A \emph{theta function} on $\mathbb W_{S,T}$ is a continuous function
  $f$ with $f(\gamma+w)=\psi(-\tfrac12\langle\gamma,w\rangle)f(w)$ for all
  matrices $\gamma$ over $E$, that is fixed by the Heisenberg action of an
  open subgroup of the finite part, and that is smooth along the archimedean
  directions.
\item A \emph{Weil datum} for $(S,T)$ consists of a continuous character
  $\psi$ of $\mathbf A_F/F$ whose archimedean component is
  $x\mapsto e^{-2\pi i\operatorname{tr}(x)}$, unitary Hecke characters
  $\chi_S,\chi_T$ of $\mathbf A_E^\times/E^\times$ with
  $\chi_S|_{\mathbf A_F^\times}=\eta^s$ and $\chi_T|_{\mathbf A_F^\times}=\eta^t$,
  and a representation $\omega$ of $U(S)(\mathbf A_F)\times U(T)(\mathbf A_F)$
  on functions on $\mathbb W_{S,T}$ that preserves theta functions, intertwines
  the Heisenberg action with the linear action of the two groups, fixes the
  value at $0$ under $U(S)(F)\times U(T)(F)$, is continuous, preserves the
  $L^2$-norm over a fundamental domain for the rational points, and acts on the
  split components through the characters $\chi_S,\chi_T$.
\item The \emph{theta kernel} of a theta function $f$ is
  $\theta(g,h;f)=(\omega(g,h)f)(0)$.
\end{enumerate}
For a Hermitian space $V$ of dimension $n$, $V^\square=V\oplus(-V)$ denotes
the doubled space, $P\subset U(V^\square)$ the Siegel parabolic subgroup
stabilising the diagonal, and $[U(T)]=U(T)(F)\backslash U(T)(\mathbf A_F)$.
\end{definition}
\leanref{Formalism/LocalToGlobal/IndefiniteReal/LocalNewformFunction.lean}{Formalism.IndefR.Datum.IsThetaFunO}
\leanref{Formalism/LocalToGlobal/IndefiniteReal/LocalNewformFunction.lean}{Formalism.IndefR.Datum.IsWeilDatumO}
\leanref{Formalism/LocalToGlobal/IndefiniteReal/LocalNewformFunction.lean}{Formalism.IndefR.Datum.thetaKO}

\begin{definition}\label{lit-siegel-eisenstein}
Let $\mathcal D$ and $E/F$ be as in Definition~\ref{lit-weil-datum}, let $V$
be a nondegenerate Hermitian space over $E$ of dimension $n\ge1$ with a chosen
basis, and let $P\subset U(V^\square)$ be the Siegel parabolic subgroup
stabilising the diagonal $\{(v,v)\}$. For $p\in P(\mathbf A_F)$ let
$a(p)\in GL_n(\mathbf A_E)$ be the matrix of its action on the diagonal, and
for $x\in\mathbf A_E^\times$ put $|x|=|N_{E/F}(x)|_{\mathbf A_F}$, with the
normalized absolute values at all places. Let $a$ be a majorant of $V$ and
let $K_a\subset U(V^\square)(\mathbf A_F)$ be the compact subgroup attached to
$a$ and the lattice of $V^\square$ spanned by the chosen basis twice. Assume
that $U(V^\square)(\mathbf A_F)=P(\mathbf A_F)K_a$.
\begin{enumerate}
\item For a homomorphism $\mu:\mathbf A_E^\times\to\mathbf C^\times$ with
  values of absolute value $1$ and for $s\in\mathbf C$, the
  \emph{degenerate principal series} $I(s,\mu)$ is the space of continuous
  functions $\Phi$ on $U(V^\square)(\mathbf A_F)$ whose right translates
  under $K_a$ span a finite-dimensional space and which satisfy
  \[
  \Phi(pg)=\mu(\det a(p))\,|\det a(p)|^{s+n/2}\,\Phi(g)
  \qquad(p\in P(\mathbf A_F),\ g\in U(V^\square)(\mathbf A_F)).
  \]
\item For $\varphi\in I(s_1,\mu)$ the \emph{flat section} through $\varphi$ is
  the family $\Phi(s)\in I(s,\mu)$, $s\in\mathbf C$, defined by
  $\Phi(s)(pk)=\varphi(pk)\,|\det a(p)|^{s-s_1}$ for $p\in P(\mathbf A_F)$
  and $k\in K_a$. It agrees with $\varphi$ on $K_a$ and $\Phi(s_1)=\varphi$.
\item The \emph{Siegel Eisenstein series} of a function $\Phi$ on
  $U(V^\square)(\mathbf A_F)$ is
  \[
  E(g,\Phi)=\sum_{\gamma\in P(F)\backslash U(V^\square)(F)}\Phi(\gamma g),
  \]
  at those $g$ where $\Phi$ is invariant under $P(F)$ on the left and the sum
  converges absolutely. For a flat section we write $E(g,\Phi(s))$.
\end{enumerate}
\end{definition}
\leanref{Development/AbelianVarietyTheory/Definitions/Foundations/T2SpaceAdeleRing.lean}{GlobalTheory.Doubling.degPS}
\leanref{Development/AbelianVarietyTheory/Definitions/Foundations/MemSiegelOfIffRightAct.lean}{GlobalTheory.Doubling.degPSOf}
\leanref{Development/AbelianVarietyTheory/Definitions/Foundations/IsIwasawaOfIsCompact.lean}{GlobalTheory.Doubling.IndLaw}
\leanref{Development/AbelianVarietyTheory/Definitions/Foundations/T2SpaceAdeleRing.lean}{GlobalTheory.Doubling.flat}
\leanref{Development/AbelianVarietyTheory/Definitions/Foundations/IsIwasawaOfIsCompact.lean}{GlobalTheory.Doubling.flatSec}
\leanref{Development/AbelianVarietyTheory/Definitions/Foundations/EisValueAtIffEisPoleLE.lean}{GlobalTheory.Eis.eis}

\begin{theorem}[Weil \cite{Wei65}]\label{lit-weil-siegel}
Let $\mathcal D$ be a definite-side datum (Definition~\ref{th-definite-datum})
and let $A,B$ be nondegenerate Hermitian spaces over its CM field $E$, with
$\operatorname{rk}A\ge1$. Let $r_B$ be the Witt index of $B$, and let $dh$ be
a Haar measure on $U(B)(\mathbf A_F)$. Assume that $B$ is anisotropic or that
$\operatorname{rk}A+r_B<\operatorname{rk}B$. Then:
\begin{enumerate}
\item for every basis of $E$ over $F$, identifying the space of
  $\operatorname{rk}B\times\operatorname{rk}A$ matrices over $\mathbf A_E$ with
  $\mathbf A_F^{r}$, $r=2\operatorname{rk}A\operatorname{rk}B$, there is a
  seminorm $\nu$ on the Schwartz space of $\mathbf A_F^{r}$, continuous in the
  archimedean component for each fixed finite component, with
  \[
  \int_{[U(B)]}\sum_{\xi}|\phi(h^{-1}\xi)|\,dh\le\nu(\phi)
  \]
  for every Schwartz function $\phi$, the sum running over the
  $\operatorname{rk}B\times\operatorname{rk}A$ matrices $\xi$ over $E$;
\item for every Weil datum $(\psi,1,\chi,\omega)$ for
  $(A^\square,B)$ (Definition~\ref{lit-weil-datum}) and every theta function
  $f$, the integral of $h\mapsto\theta(g,h;f)$ over $[U(B)]$ converges
  absolutely for every $g$;
\item if moreover $2\operatorname{rk}A<\operatorname{rk}B$, then for every Haar
  measure of covolume one on $\operatorname{Hom}(A,B)(\mathbf A_F)$ with
  fundamental domain $\mathcal G$ for the rational points, the Siegel-Weil
  section $\Phi_f(g)=\int_{\mathcal G}(\omega(g,1)f)(y,-y)\,dy$, where
  $(y,-y)$ is the $\operatorname{rk}B\times2\operatorname{rk}A$ matrix with
  blocks $y$ and $-y$, is invariant
  under $P(F)$, the series $\sum_{\gamma\in P(F)\backslash U(A^\square)(F)}
  \Phi_f(\gamma g)$ converges absolutely, and its value is the average of
  $\theta(g,\cdot\,;f)$ over $[U(B)]$ for the probability measure.
\end{enumerate}
\end{theorem}
\leanref{Development/AbelianVarietyTheory/PublishedTheorems/TwistedTraceFormula/WeilSiegelFormulaUnitary.lean}{GlobalTheory.Published.WeilSiegelFormulaUnitary}
\leanref{Development/AutomorphicArithmetic/ThetaCorrespondence/FockModels/FockAnisoFree.lean}{UpperTheory.Lit.Sw47.weilUnitary}

\begin{theorem}[Siegel-Weil formula in the convergent range \cite{KR88,Ich04,Ich07}]\label{lit-siegel-weil}
Let $\mathcal D$ be a definite-side datum (Definition~\ref{th-definite-datum})
and let $A,B$ be nondegenerate Hermitian spaces over its CM field $E$, with
$\operatorname{rk}A<\operatorname{rk}B\le2\operatorname{rk}A$ and
$\operatorname{rk}A+r_B<\operatorname{rk}B$, where $r_B$ is the Witt index of
$B$. Let $(\psi,1,\chi_B,\omega)$ be a Weil datum for $(A^\square,B)$
(Definition~\ref{lit-weil-datum}). Let $a$ be a majorant of $A$ as in
Definition~\ref{lit-siegel-eisenstein}. Choose archimedean data for
$\mathcal D$ and Cartan involutions $s_A$ of $A^\square$ and $s_B$ of $B$ at
the archimedean places such that $s_A$ commutes with the archimedean part of
$K_a$. Fix a Haar measure of covolume one on
$\operatorname{Hom}(A,B)(\mathbf A_F)$ and a Haar measure on
$U(B)(\mathbf A_F)$. Let $f$ be a Fock vector: a theta function annihilated
by every product of $N$ archimedean annihilation operators attached to
$(s_A,s_B)$, for some $N$. Put $s_0=(\operatorname{rk}B-\operatorname{rk}A)/2$
and let $\Phi(s)$ be the flat section through the Siegel-Weil section
$\Phi_f$ of Theorem~\ref{lit-weil-siegel}(3) at $s_0$. Then the integral
of $h\mapsto\theta(g,h;f)$ over $[U(B)]$ converges absolutely for every $g$,
and for every $g$ the function $s\mapsto E(g,\Phi(s))$ continues analytically
to a neighbourhood of $s_0$ with value at $s_0$
\[
E(g,\Phi(s_0))=\int_{[U(B)]}\theta(g,h;f)\,dh,
\]
the integral taken for the probability measure.
\end{theorem}
\leanref{Development/AbelianVarietyTheory/PublishedTheorems/SiegelWeil₂/SiegelWeilUnitaryConvergent.lean}{GlobalTheory.Published.SiegelWeilUnitaryConvergent}
\leanref{Development/AutomorphicArithmetic/ThetaCorrespondence/FockModels/FockAnisoFree.lean}{UpperTheory.Lit.Sw47.convergent}

\begin{theorem}[Doubling method \cite{GPSR87}]\label{lit-doubling}
Let $\mathcal D$, $V$, $n$, $a$ and $K_a$ be as in
Definition~\ref{lit-siegel-eisenstein}. Let $\mu$ be a Hecke character of
$\mathbf A_E^\times/E^\times$ with $\mu|_{\mathbf A_F^\times}=\eta^b$ for some
$b\ge0$, let $s_1\in\mathbf C$, let $\varphi\in I(s_1,\mu)$, and let
$\Phi(s)$ be the flat section through $\varphi$. Then:
\begin{enumerate}
\item for $\operatorname{Re}s>n/2$ the section $\Phi(s)$ is invariant under
  $P(F)$ on the left and the series $E(g,\Phi(s))$ converges absolutely for
  every $g$; for every $g$ the function $s\mapsto E(g,\Phi(s))$ continues
  meromorphically to $\mathbf C$; and for every $s_0$ there are $k\ge0$ and
  $r>0$ such that $(s-s_0)^kE(g,\Phi(s))$ extends to a function on the disc
  $|s-s_0|<r$ that is holomorphic in $s$, jointly continuous in $(s,g)$, and
  bounded by $C\,h(g)^N$ for some $C,N$ independent of $s$, where $h$ is
  the height function on $U(V^\square)(\mathbf A_F)$;
\item assume that the signature $(p_\tau,q_\tau)$ recorded with $V$ at each
  archimedean place $\tau$ is that of its Hermitian form, choose archimedean
  data for $\mathcal D$ and for each $\tau$ a matrix $A_\tau$ with
  $A_\tau^*J_{p_\tau,q_\tau}A_\tau$ equal to the Gram matrix of $V$ at $\tau$,
  where $J_{p,q}=\operatorname{diag}(1_p,-1_q)$, and let $f_1,f_2$ be cusp
  forms on $U(V)$ that are finite under the maximal compact subgroup at
  infinity determined by the $A_\tau$. Let $m$ be a Haar measure on $U(V)(\mathbf A_F)$. Then for
  $\operatorname{Re}s>n/2$ the doubling integral
  \[
  Z(s)=\int_{[U(V)]\times[U(V)]}E\bigl((g_1,g_2),\Phi(s)\bigr)\,
  \overline{f_1(g_1)}\,f_2(g_2)\,\mu(\det g_2)^{-1}\,dg_1\,dg_2
  \]
  converges absolutely, the integral
  $\int_{U(V)(\mathbf A_F)}\Phi(s)\bigl((g,1)\bigr)\langle f_1(\cdot\,g),f_2\rangle\,dg$
  converges absolutely, and the two are equal; and $Z(s)$ continues
  meromorphically to $\mathbf C$.
\end{enumerate}
\end{theorem}
\leanref{Development/WeilBound/RepresentationTheoryGL/EisensteinResidues/AnisotropicGlobalForms/DmOfHyp.lean}{ResidueKilled.DM}
\leanref{Development/AutomorphicArithmetic/ThetaCorrespondence/ThetaPairLifts/LocalTypeDictionary/DiscreteSpectrumResidualClassification.lean}{GlobalTheory.Published.DoublingMethodLocalZetaIntegrals}
\leanref{Development/WeilBound/RepresentationTheoryGL/EisensteinResidues/AnisotropicGlobalForms/DmOfHyp.lean}{ResidueKilled.dm_of_two}
\leanref{Development/AutomorphicArithmetic/ThetaCorrespondence/ThetaPairLifts/ThirdPlaceSeries/EveryCharacterOfLiteratureOfVWin.lean}{UpperTheory.LiteratureTheorems.doubling_method}

\begin{remark}\label{lit-vacuous-slot}
In the setting of Theorem~\ref{lit-doubling}, Alpöge also states a form of
the doubling method with a third assertion: when $n<b$, the Eisenstein series
$E(g,\Phi(s))$ has at most a simple pole at $s=(b-n)/2$, as in \cite{Tan99}.
The list of Theorem~\ref{lit-bundle-unitary} has an entry named for this
simple pole, and its content is the true proposition. The doubling entry of
the same list is Theorem~\ref{lit-doubling}, in which the pole assertion is
replaced by the true proposition as well. So the simple-pole statement is not
proved, and the proof of the main theorems does not use it. Theorem~\ref{lit-doubling} is proved from rapid decay of
cusp forms and the meromorphic continuation of Siegel Eisenstein series.
\end{remark}
\leanref{Development/AutomorphicArithmetic/ThetaCorrespondence/ThetaPairLifts/LocalTypeDictionary/DiscreteSpectrumResidualClassification.lean}{GlobalTheory.Published.UnusedSlot72}
\leanref{Development/AutomorphicArithmetic/ThetaCorrespondence/ThetaPairLifts/ThirdPlaceSeries/EveryCharacterOfLiteratureOfVWin.lean}{UpperTheory.LiteratureTheorems.siegel_eisenstein_simple_pole}
\leanref{Development/AbelianVarietyTheory/PublishedTheorems/SiegelWeil₂/SiegelEisensteinSimplePoleOfDoublingMethod.lean}{GlobalTheory.Published.SiegelEisensteinSimplePole}
\leanref{Development/AbelianVarietyTheory/PublishedTheorems/TwistedTraceFormula/EisRegularAtOfEisPoleLE.lean}{GlobalTheory.Published.DoublingMethod}

\section{Automorphic representations of general linear groups}\label{section-lit-gln}

\begin{theorem}[Jacquet-Shalika \cite{JS81a,JS81b}]\label{lit-jacquet-shalika}
Let $E$ be a totally complex number field.
\begin{enumerate}
\item Let $(\pi_i)_{i\in I}$ and $(\pi'_j)_{j\in J}$ be finite families of
  cuspidal automorphic representations of $GL_{m_i}(\mathbf A_E)$ and
  $GL_{m'_j}(\mathbf A_E)$, with all $m_i,m'_j\ge1$. Assume that for all but
  finitely many finite places $w$ of $E$ every $\pi_i$ and every $\pi'_j$ has
  Satake parameters at $w$, and the sum of the multisets of Satake parameters
  of the $\pi_{i,w}$ equals that of the $\pi'_{j,w}$. Then there is a bijection
  $s:I\to J$ with $m'_{s(i)}=m_i$ and $\pi'_{s(i)}\cong\pi_i$ for all $i$.
\item Let $\pi$ be a cuspidal automorphic representation of
  $GL_a(\mathbf A_E)$ with unitary central character, and let $w$ be a finite
  place at which $\pi$ has Satake parameters, with $q_w$ the norm of $w$.
  Then every Satake parameter $\alpha$ of $\pi_w$ satisfies
  $q_w^{-1/2}<|\alpha|<q_w^{1/2}$.
\end{enumerate}
\end{theorem}
\leanref{Development/AbelianVarietyTheory/PublishedTheorems/GaloisCohomology₂/StrongMultiplicityOne.lean}{GlobalTheory.Published.JacquetShalika}
\leanref{Development/WeilBound/CoreStatements/PublishedTheorems/JacquetShalikaRigidity/JacquetShalikaRigidity.lean}{PositivityTheory.AdelicQuotientData.jacquetShalika}

\begin{theorem}[Clozel \cite{Clo90}]\label{lit-clozel-cohomological}
Let $E$ be a totally complex number field, let $a\ge1$, and let $\pi$ be an
irreducible cuspidal automorphic representation of $GL_a(\mathbf A_E)$. Assume
that $\pi$ is regular algebraic: for every embedding $\tau:E\to\mathbf C$ the
infinitesimal character of $\pi$ at $\tau$ consists of $a$ distinct numbers of
the form $m+\frac{a-1}2$ with $m\in\mathbf Z$. Then $\pi$ is cohomological:
there is an irreducible algebraic representation $\Xi$ of
$GL_a(E\otimes_{\mathbf Q}\mathbf C)$ such that $\pi_\infty\otimes\Xi$ has
nonzero $(\mathfrak g,K_\infty)$-cohomology.
\end{theorem}
\leanref{Development/AbelianVarietyTheory/PublishedTheorems/GaloisGeometry₂/ClozelCohomologicalModel.lean}{GlobalTheory.Published.ClozelCohomological}
\leanref{Development/WeilBound/RepresentationTheoryGL/CasimirAlgebra/KugaCoordinates/ClozelCohomological.lean}{PositivityTheory.Rg.Cz.clozelCohomological}

\section{$p$-adic Hodge theory}\label{section-lit-hodge}

\begin{theorem}[Bhatt-Scholze \cite{BS23}]\label{lit-bhatt-scholze}
Let $\ell$ be a prime, let $K$ be a finite extension of $\mathbf Q_\ell$ with
ring of integers $\mathcal O_K$, and identify the $A_{\inf}$-prism of
$\mathcal O_K$ with $A_{\inf}$, compatibly with Frobenius, the ideal
$\ker\theta$ and the map to $\mathcal O_{\mathbf C_\ell}$. For a prismatic
$F$-crystal $\mathcal E$ in vector bundles on $\mathcal O_K$ and a
$\mathbf Z_\ell$-module $T$ with an action of $G_K$, call a semilinear map
$j:\mathcal E(A_{\inf})\to W(\mathbf C_\ell^\flat)\otimes_{\mathbf Z_\ell}T$ a
realisation if it becomes an isomorphism after base change along
$A_{\inf}\to W(\mathbf C_\ell^\flat)$ and is compatible with Frobenius and
with the action of $G_K$. Then:
\begin{enumerate}
\item for every finite extension $L$ of $\mathbf Q_\ell$ and every
  $G_K$-stable $\mathcal O_L$-lattice $T$ in a crystalline representation of
  $G_K$ on a finite-dimensional $L$-vector space, there is a prismatic
  $F$-crystal $\mathcal E$ in vector bundles with a realisation
  $\mathcal E(A_{\inf})\to W(\mathbf C_\ell^\flat)\otimes T$;
\item for realisations $j:\mathcal E(A_{\inf})\to W(\mathbf C_\ell^\flat)\otimes T$
  and $j':\mathcal E'(A_{\inf})\to W(\mathbf C_\ell^\flat)\otimes T'$, every
  $G_K$-equivariant $\mathbf Z_\ell$-linear map $f:T\to T'$ is induced by a
  unique map of $F$-crystals $u:\mathcal E\to\mathcal E'$, in the sense that
  $j'\circ u=(1\otimes f)\circ j$.
\end{enumerate}
\end{theorem}
\leanref{Development/AbelianVarietyTheory/PublishedTheorems/GaloisCohomology₂/BhattScholzeCrystallineLattices.lean}{GlobalTheory.Published.BhattScholzeCrystallineLattices}
\leanref{Development/AutomorphicArithmetic/ThetaCorrespondence/ThetaPairLifts/ThirdPlaceSeries/EveryCharacterOfLiteratureOfVWin.lean}{UpperTheory.LiteratureTheorems.bhatt_scholze_crystalline_lattices_of_literature}

\begin{theorem}[Fargues \cite{Far08}]\label{lit-hodge-tate-sequence}
Let $p$ be a prime, let $C$ be an algebraically closed field of
characteristic zero, complete for a rank-one valuation with $|p|<1$, and let
$\mathcal O$ be its valuation ring. Let $G$ be a $p$-divisible group over
$\mathcal O$ with dual $G^\vee$, Tate modules $T_pG$ and $T_pG^\vee$, and
Hodge-Tate maps
$\alpha_G:\mathcal O\otimes T_pG\to\omega_{G^\vee}$ and
$\alpha_{G^\vee}:\mathcal O\otimes T_pG^\vee\to\omega_G$. Fix a compatible
system of primitive $p$-power roots of unity, and let $B$ be the
$\mathcal O$-bilinear form on $(\mathcal O\otimes T_pG)\times
(\mathcal O\otimes T_pG^\vee)$ induced by the Weil pairing. Then for every
$x\in\mathcal O\otimes T_pG$,
\[
\alpha_G(x)=0\iff B(x,y)=0\text{ for every }y\in\ker\alpha_{G^\vee}.
\]
\end{theorem}
\leanref{Development/AutomorphicSpectrum/PublishedTheorems/Algebra₂/FarguesExact.lean}{GlobalTheory.CompanionTheory.Ar.HT.FarguesExact}
\leanref{Development/AutomorphicArithmetic/ThetaCorrespondence/PrismaticDieudonne/PrismaticDieudonneHodge.lean}{UpperTheory.Lit.farguesExact}

\section{Analytic and algebraic number theory}\label{section-lit-number-theory}

\begin{theorem}[Linnik \cite{Lin44}]\label{lit-linnik}
There are constants $C,L>0$ such that for all integers $a,q\ge1$ with
$\gcd(a,q)=1$ the least prime $p\equiv a\pmod q$ satisfies $p\le Cq^L$.
\end{theorem}
\leanref{Development/AbelianVarietyTheory/PublishedTheorems/PublishedFoundations/PrimesInProgressionNonempty.lean}{GlobalTheory.Published.Linnik}
\leanref{Development/AbelianVarietyTheory/PublishedTheorems/PublishedFoundations/PrimesInProgressionNonempty.lean}{GlobalTheory.Published.linnik_iff_exists_prime}
\leanref{Development/AutomorphicArithmetic/ThetaCorrespondence/Linnik.lean}{UpperTheory.Lit.linnik}

\begin{theorem}[Chebotarev \cite{Che26,Neu99}]\label{lit-chebotarev-degree-one}
Let $F$ be a number field, let $L/F$ be a finite Galois extension, let
$g\in\operatorname{Gal}(L/F)$, and let $S$ be a finite set of finite places
of $F$. Then there is a finite place $v\notin S$ of $F$ of absolute residue
degree one, unramified in $L$, with a Frobenius element in
$\operatorname{Gal}(\bar F/F)$ whose restriction to $L$ is $g$.
\end{theorem}
\leanref{Development/AbelianVarietyTheory/PublishedTheorems/GaloisCohomology₂/ChebotarevDegreeOneIffCheb.lean}{GlobalTheory.Published.ChebotarevDegreeOne}
\leanref{Development/AbelianVarietyTheory/Proofs/InertiaEigenvalues/ChebotarevDegreeOne.lean}{GlobalTheory.Published.chebotarevDegreeOne'}

\section{Resolution and arithmetic geometry}\label{section-lit-geometry}

The covolume bound of Chapter~\ref{chapter-covolumes} uses resolution of
singularities, principalization, an arithmetic Hilbert-Samuel estimate and
two results on abelian schemes and $p$-divisible groups. A closed subset of a
scheme is \emph{simple normal crossings} at a point $x$ if the local ring at
$x$ is regular and, on an affine neighbourhood, the vanishing ideal of the
subset is generated by a product of some members of a family of functions
whose germs form a regular system of parameters at $x$.

\begin{definition}\label{lit-snc-sequence}
Let $k$ be a field, let $X$ be a smooth $k$-scheme and let $E\subset X$ be a
closed subset. A \emph{normal crossings blow-up sequence} starting at $(X,E)$
is a finite sequence $X_r\to\cdots\to X_0=X$ in which each $X_{i+1}\to X_i$ is
the blow-up of a closed subscheme $Z_i\subset X_i$, together with closed
subsets $E_0=E$ and $E_{i+1}=$ the inverse image of $E_i\cup Z_i$, such that
each $X_i$ and each $Z_i$ is smooth over $k$, and $Z_i$ has normal crossings
with $E_i$: at every point of $Z_i$ there is a regular system of parameters
such that $E_i$ is cut out by a product of some of them and $Z_i$ by some of
them.
\end{definition}
\leanref{Development/AutomorphicSpectrum/PublishedTheorems/Geometry₂/NotSubsetOfDenseCompl.lean}{GlobalTheory.CompanionTheory.Hypotheses₂.SncBlowup.IsSncBlowupSeq}
\leanref{ReciprocityTheory/AnalyticForms/RegularLocalRings/FinalTheorems.lean}{ReciprocityTheory.AnForm.R.FinalTheorems.Snc}

\begin{theorem}[Hironaka \cite{Hir64,Kol07}]\label{lit-principalization}
Let $k$ be a field of characteristic zero and let $X$ be a quasi-compact
separated scheme, smooth of relative dimension $d$ over $k$. Let
$I\subset\mathcal O_X$ be an ideal sheaf whose zero set contains no
irreducible component of $X$, and let $E\subset X$ be a closed subset that is
simple normal crossings at each of its points. Then there is a normal
crossings blow-up sequence $s:X'\to X$ starting at $(X,E)$
(Definition~\ref{lit-snc-sequence}) such that $X'$ is smooth over $k$, the
ideal $I\mathcal O_{X'}$ is invertible, the reduced zero set of
$I\mathcal O_{X'}$ is a simple normal crossings divisor, and $s$ is an
isomorphism over the complement of the zero set of $I$. For $E=\emptyset$ compare
\cite[Tag R021]{Res26}.
\end{theorem}
\leanref{Development/AutomorphicSpectrum/PublishedTheorems/Geometry₂/NotSubsetOfDenseCompl.lean}{GlobalTheory.CompanionTheory.Hypotheses₂.PrincipalisationCharZero}
\leanref{Development/WeilBound/CoreStatements/ResolutionCharZero/PrincipalisationCharZeroClosed.lean}{PositivityTheory.ResPlug.WeightsCompactStabilizer.principalisationCharZero_closed}

\begin{corollary}\label{lit-principalization-projective}
Let $k$ be a field of characteristic zero, let $Y$ be an integral scheme,
smooth of relative dimension $d$ and projective over $k$, and let
$I\subset\mathcal O_Y$ be a nonzero ideal sheaf. Then there is a finite
sequence of blow-ups $s:Y'\to Y$ in centres smooth over $k$ lying over the
zero set of $I$, such that $Y'$ is integral, smooth of relative dimension $d$
and projective over $k$, $s$ is an isomorphism over the complement of the zero
set of $I$, the ideal $I\mathcal O_{Y'}$ is invertible, and its reduced zero
set is a simple normal crossings divisor.
\end{corollary}
\begin{proof}
Since $Y$ is integral and $I\neq0$, the zero set of $I$ is a proper closed
subset and contains no irreducible component. Apply
Theorem~\ref{lit-principalization} with $E=\emptyset$, and write
$Y'=X_r\to\cdots\to X_0=Y$ for the sequence, with centres $Z_i$ smooth over
$k$. Projectivity over $k$ is preserved by each blow-up. The complement $U$ of
the zero set of $I$ is nonempty and $s$ is an isomorphism over it, so $Y'$ is
nonempty. Induction along the sequence, using that each $X_i$ is smooth and
$Y'$ is nonempty, shows that every $X_i$ is integral. The inverse image of $U$
is dense in $Y'$, because the complement of a simple normal crossings divisor
is dense. Over it $Y'$ has relative dimension $d$, and $Y'$ is smooth and
irreducible, so it has relative dimension $d$ everywhere.

It remains to move the centres over the zero set of $I$. Every blow-up
$X_{i+1}\to X_i$ is an isomorphism over the inverse image $U_i$ of $U$. On the
smooth integral scheme $X_i$, the components of $Z_i$ of codimension one form
an open and closed part of $Z_i$ along which $Z_i$ is a Cartier divisor.
Removing them leaves a centre $Z'_i$, smooth over $k$, with the same blow-up,
and $Z'_i$ is disjoint from $U_i$ because a blow-up in a smooth centre of
codimension at least two is not an isomorphism over any point of the centre.
\end{proof}
\leanref{Development/AbelianVarietyTheory/PublishedTheorems/Divisors₂/GeometricallyConnectedOfIsIsoRestrict.lean}{GlobalTheory.Published.Principalisation}
\leanref{Development/WeilBound/CoreStatements/ResolutionCharZero/PrincipalisationClosed.lean}{PositivityTheory.ResPlug.WeightsCompactStabilizer.principalisation_closed}

\begin{corollary}\label{lit-resolution}
Let $k$ be a field of characteristic zero and let $X$ be a reduced scheme,
separated, quasi-compact and locally of finite type over $k$, with smooth locus
$X_{\mathrm{sm}}$. Then there is a finite sequence of blow-ups $s:X'\to X$ such
that $X'$ is smooth over $k$, $s$ is an isomorphism over $X_{\mathrm{sm}}$,
and the reduced inverse image of $X\setminus X_{\mathrm{sm}}$ is a simple
normal crossings divisor. This is a weak form of \cite[Main Theorem I]{Hir64}
and \cite[Theorem 36]{Kol07}; compare \cite[Tag R02G]{Res26}.
\end{corollary}
\begin{proof}
Theorem~\ref{cv-functorial-resolution} assigns to $X$ a finite sequence of
blow-ups with these three properties. The corollary keeps them and forgets
the compatibility of the assignment with smooth morphisms and field
extensions.
\end{proof}
\leanref{Development/AutomorphicSpectrum/PublishedTheorems/Geometry₂/ConclusionOfSmooth.lean}{GlobalTheory.CompanionTheory.Hypotheses₂.ResolutionCharZero}
\leanref{Development/WeilBound/CoreStatements/ResolutionCharZero/ResolutionCharZeroOfTop.lean}{PositivityTheory.Resplug.ResolutionBlowup.resolutionCharZero_of_top}
\leanref{Development/WeilBound/CoreStatements/ResolutionCharZero/ResolutionCharZero.lean}{PositivityTheory.Resplug.ResolutionBlowup.resolutionCharZero}

\begin{definition}\label{lit-arithmetic-variety}
An \emph{arithmetic variety} of dimension $D+1$ is an integral scheme $X$ with
a projective flat morphism $X\to\operatorname{Spec}R$, where $R$ is a domain
finite over $\mathbf Z$, such that $X$ has Krull dimension $D+1$. A
\emph{Hermitian line bundle} $\bar L$ on $X$ is a line bundle with a smooth
metric on $L_\sigma$ over $X_\sigma(\mathbf C)$ for each embedding
$\sigma:R\to\mathbf C$, compatible with complex conjugation. It is \emph{nef}
if its metrics are semipositive, its degree on every integral curve in a
fibre over a closed point is nonnegative, and its arithmetic degree on every
horizontal integral curve is nonnegative. A global section is \emph{small} if
its pointwise norm is at most $1$ at every complex point.
\end{definition}
\leanref{Development/AbelianVarietyTheory/Definitions/Divisors/ExistsIsClosedImmersionGenericPoint.lean}{GlobalTheory.IsArithVar}
\leanref{Development/AbelianVarietyTheory/Definitions/Divisors/ExistsIsClosedImmersionGenericPoint.lean}{GlobalTheory.HermLB}
\leanref{Development/AbelianVarietyTheory/Definitions/Divisors/ExistsIsClosedImmersionGenericPoint.lean}{GlobalTheory.HermLB.IsNef}
\leanref{Development/AbelianVarietyTheory/Definitions/Divisors/ExistsIsClosedImmersionGenericPoint.lean}{GlobalTheory.HermLB.smallSec}

\begin{theorem}[Arithmetic Hilbert-Samuel formula \cite{GS92,Zha95,Mor09}]\label{lit-arithmetic-hilbert-samuel}
Let $X$ be an arithmetic variety of dimension $D+1$
(Definition~\ref{lit-arithmetic-variety}), let $\bar L$ and $\bar N$ be
Hermitian line bundles on $X$ with $\bar L$ nef, and let $a$ be the arithmetic
self-intersection number of $\bar L$. Then for all large $m$ the set of small
sections of $\bar L^{\otimes m}\otimes\bar N$ is finite, and
\[
\log\#\{\text{small sections of }\bar L^{\otimes m}\otimes\bar N\}
=\frac{a}{(D+1)!}\,m^{D+1}+o(m^{D+1})\qquad(m\to\infty).
\]
\end{theorem}
\leanref{Development/AbelianVarietyTheory/PublishedTheorems/Divisors₂/SmallSecTensTriv.lean}{GlobalTheory.Published.ArithHilbertSamuel}
\leanref{Development/AbelianVarietyTheory/Proofs/DivisorProofs/ArithHilbertSamuelOfPrincipalisation.lean}{GlobalTheory.arithHilbertSamuel_of_principalisation}

\begin{definition}\label{lit-multiprojective}
Let $K$ be a finite nonempty set, for each $k\in K$ let $I_k$ be a finite
nonempty set, and let $R=\mathbf Z[x_{k,i}:k\in K,\ i\in I_k]$, graded by
multidegree in $\mathbf Z^K$. Let $\hat C\subset\prod_k\mathbf C^{I_k}$ be the
set of points $z=(z_{k,i})$ with $(z_{k,i})_{i\in I_k}\neq0$ for every $k$.
\begin{enumerate}
\item A multihomogeneous prime ideal $W\subset R$ is \emph{relevant} if for
  each $k$ some $x_{k,i}$ is not in $W$. An \emph{arithmetic subvariety of
  dimension $d$} is a relevant multihomogeneous prime $P$ with
  $P\cap\mathbf Z=0$ and $\dim R/P=d+|K|$; it is an integral closed subscheme
  $X$ of $\prod_k\mathbf P^{I_k}_{\mathbf Z}$, flat over $\mathbf Z$, of
  dimension $d$, and we write $\hat X(\mathbf C)$ for the set of
  $z\in\hat C$ at which every element of $P$ vanishes. A \emph{horizontal
  curve} in $X$ is an arithmetic subvariety of dimension $1$ containing $P$; a
  \emph{vertical curve} in $X$ is a relevant multihomogeneous prime
  $W\supset P$ with $W\cap\mathbf Z\neq0$ and $\dim R/W=1+|K|$.
\item A \emph{metrized line bundle} $\bar A=(a,h)$ consists of
  $a\in\mathbf Z^K$ and a positive function $h$ on $\hat C$ with
  $h(cz)=\prod_k|c_k|^{a_k}h(z)$ for $c\in(\mathbf C^\times)^K$ acting
  blockwise. It is a metric on $\mathcal O(a)$: a multihomogeneous $m$ of
  degree $a$ has norm $|m(z)|/h(z)$ at the point of $z$. The bundle $\bar A$ is
  \emph{smooth} if $h$ is smooth, \emph{invariant under conjugation} if
  $h(\bar z)=h(z)$, \emph{semipositive on $X$} if it is smooth and the Levi
  form of $\log h$ at each $z\in\hat X(\mathbf C)$ is nonnegative on the
  vectors $w$ with $\sum_{k,i}\frac{\partial f}{\partial x_{k,i}}(z)w_{k,i}=0$
  for all $f\in P$, and \emph{nef on the fibres of $X$} if $\mathcal O(a)$ has
  nonnegative degree on every vertical curve in $X$.
\item An \emph{arithmetic intersection system} assigns to every such $K$, every
  arithmetic subvariety $X$ of dimension $n$ and all smooth metrized line
  bundles $\bar L_1,\ldots,\bar L_n$ a real number
  $(\bar L_1\cdots\bar L_n\mid X)$ that is symmetric in the $\bar L_j$,
  additive in each of them, and satisfies the following recursion: if
  $s\in R\setminus P$ is multihomogeneous of degree $a$, where $\bar A=(a,h)$,
  then $(\bar L_1\cdots\bar L_{n-1}\bar A\mid X)$ is the sum over the relevant
  minimal primes $W$ of $P+(s)$, with the multiplicity of $W$ in $R/(P+(s))$,
  of $(\bar L_1\cdots\bar L_{n-1}\mid W)$ if $W$ is horizontal and of
  $\deg_{L_1\cdots L_{n-1}}(W)\log p$ if $W\cap\mathbf Z=p\mathbf Z$, minus the
  integral of $\log(|s|/h)$ over $\hat X(\mathbf C)$ against the product of
  the first Chern forms of $\bar L_1,\ldots,\bar L_{n-1}$.
\item Let $\Omega$ be a field and $y\in\Omega^{\sqcup_kI_k}$, and let $P$ be the
  kernel of $R\to\Omega[t_k:k\in K]$, $x_{k,i}\mapsto y_{k,i}t_k$. The field
  $\Omega$ is the \emph{function field} of $P$ if every element of $\Omega$ is
  $g(y)/h(y)$ with $g,h$ multihomogeneous of the same degree. For $\tau\ge0$ let
  $H^0(X,\tau\bar A)$ be the group of $s\in\Omega$ such that
  $s\prod_ky_{k,c_k}^{-\tau a_k}$ lies in
  $\mathbf Z[y_{k,i}/y_{k,c_k}:k\in K,\ i\in I_k]$ for every choice of indices
  $c_k\in I_k$ with all $y_{k,c_k}\neq0$, and let $\tilde H^0(X,\tau\bar A)$
  be defined in the same way with the integral closures of these rings in
  $\Omega$. An element $s\in\Omega$ is \emph{small} for $\tau\bar A$ if
  $|g(z)|\le|h(z)|\,h_A(z)^\tau$ for all $z\in\hat X(\mathbf C)$ and all
  multihomogeneous $g,h$ with $\deg g=\deg h+\tau a$ and $g(y)=s\,h(y)$, where
  $h_A$ is the function of $\bar A$.
\end{enumerate}
\end{definition}
\leanref{ReciprocityTheory/ReductiveGroups/ArithmeticGeometry/BirationalGeometry/IsIntersectionFamily.lean}{ReciprocityTheory.ChartRestriction.G.Geo.IsArithVar}
\leanref{ReciprocityTheory/ReductiveGroups/ArithmeticGeometry/BirationalGeometry/IsIntersectionFamily.lean}{ReciprocityTheory.ChartRestriction.G.Geo.HermLine}
\leanref{ReciprocityTheory/ReductiveGroups/ArithmeticGeometry/BirationalGeometry/ArithmeticHilbertSamuelLowerBound.lean}{ReciprocityTheory.ChartRestriction.G.Geo.IsSemipositiveRel}
\leanref{ReciprocityTheory/ReductiveGroups/ArithmeticGeometry/BirationalGeometry/ArithmeticHilbertSamuelLowerBound.lean}{ReciprocityTheory.ChartRestriction.G.Geo.NefOnFibres}
\leanref{ReciprocityTheory/ReductiveGroups/ArithmeticGeometry/BirationalGeometry/IsIntersectionFamily.lean}{ReciprocityTheory.ChartRestriction.G.Geo.IsIntersectionFamily}
\leanref{ReciprocityTheory/ReductiveGroups/ArithmeticGeometry/BirationalGeometry/IsIntersectionFamily.lean}{ReciprocityTheory.ChartRestriction.G.Geo.recR}
\leanref{ReciprocityTheory/ReductiveGroups/ArithmeticGeometry/BirationalGeometry/ArithmeticHilbertSamuelLowerBound.lean}{ReciprocityTheory.ChartRestriction.G.Geo.chartLat}
\leanref{ReciprocityTheory/ReductiveGroups/ArithmeticGeometry/BirationalGeometry/ArithmeticHilbertSamuelLowerBound.lean}{ReciprocityTheory.ChartRestriction.G.Geo.chartLatNm}
\leanref{ReciprocityTheory/ReductiveGroups/ArithmeticGeometry/BirationalGeometry/ArithmeticHilbertSamuelLowerBound.lean}{ReciprocityTheory.ChartRestriction.G.Geo.IsSmallH}

\begin{proposition}[Zhang, Moriwaki \cite{Zha95,Mor09}]\label{lit-hilbert-samuel-upper}
Let $K$, $(I_k)$ and $R$ be as in Definition~\ref{lit-multiprojective}, let
$(\cdot\mid\cdot)$ be an arithmetic intersection system, let $D\ge0$, let $X$
be an arithmetic subvariety of dimension $D+1$ defined by the prime $P$, and
let $\bar A=(a,h_A)$ be a metrized line bundle that is semipositive on $X$ and
invariant under conjugation. Write $\bar A^{D+1}$ for
$(\bar A\cdots\bar A\mid X)$.
\begin{enumerate}
\item If $\bar A$ is nef on the fibres of $X$ and $(\bar A\mid W)\ge0$ for
  every horizontal curve $W$ in $X$, then $\bar A^{D+1}\ge0$.
\item Let $\Omega$ be a field of characteristic zero and $y\in\Omega^{\sqcup_kI_k}$
  with kernel $P$ and function field $\Omega$. Assume that $D\ge1$, that
  $\mathcal O(a)$ has positive degree on the generic fibre $X_{\mathbf Q}$, and
  that there are finitely many small elements $w_1,\ldots,w_q$ of
  $H^0(X,\bar A)$ that generate $\mathcal O(a)$ on $X$: on each chart of
  Definition~\ref{lit-multiprojective}(4) the elements
  $w_j\prod_ky_{k,c_k}^{-a_k}$ generate the unit ideal. Then
  \[
  \limsup_{\tau\to\infty}\frac{\log\#\{s\in H^0(X,\tau\bar A):s\text{ small for }\tau\bar A\}}
  {\tau^{D+1}/(D+1)!}\le\bar A^{D+1}.
  \]
\item Let $\Omega$ be a field and $y\in\Omega^{\sqcup_kI_k}$ with kernel $P$ and
  function field $\Omega$. Assume that $D\ge1$, that $a_k\ge0$ for all $k$, and
  that $|m(z)|\le h_A(z)$ for every $z\in\hat X(\mathbf C)$ and every monomial
  $m$ of multidegree $a$. Then
  \[
  \limsup_{\tau\to\infty}\frac{\log\#\{s\in\tilde H^0(X,\tau\bar A):s\text{ small for }\tau\bar A\}}
  {\tau^{D+1}/(D+1)!}\le\bar A^{D+1}.
  \]
\end{enumerate}
\end{proposition}
\leanref{Development/WeilBound/CoreStatements/AutomorphicKernels/MoriwakiBigness/HilbertSamuelApplications/SmallGeneratedCountUpperOfUse.lean}{PositivityTheory.MbHsUse.HSUse}
\leanref{GroundTheory/HeightsAndTheta/Heights/HilbertSamuel/StatementsArithKleiman.lean}{GroundTheory.Heights₂.Heights.HilbertSamuel.ArithKleiman.Statement}
\leanref{GroundTheory/HeightsAndTheta/Heights/HilbertSamuel/StatementsArith.lean}{GroundTheory.Heights₂.Heights.HilbertSamuel.G7.Statement}
\leanref{ReciprocityTheory/ReductiveGroups/ArithmeticGeometry/BirationalGeometry/ArithmeticHilbertSamuelLowerBound.lean}{ReciprocityTheory.ChartRestriction.G.Geo.ArithmeticHilbertSamuelUpperBound}
\leanref{Development/WeilBound/CoreStatements/AutomorphicKernels/MoriwakiBigness/MoriwakiBigTheorem/Outright.lean}{PositivityTheory.Pf1.HilbertSamuelMoriwaki.hsUse_outright}

\begin{theorem}[Zarhin \cite{Zar74,MB85}]\label{lit-zarhin}
Let $C$ be an algebraically closed field with a valuation, let $R$ be its
valuation ring, and let $A$ be an abelian scheme over $R$ with a relatively
ample line bundle. Then there are a principally polarized abelian scheme $P$
over $R$ and homomorphisms $i:A\to P$ and $r:P\to A$ with $r\circ i=1_A$.
\end{theorem}
\leanref{Development/AutomorphicSpectrum/PublishedTheorems/Algebra₂/ZarhinStatement.lean}{GlobalTheory.CompanionTheory.Ar.GR.ZarhinStatement}
\leanref{Development/WeilBound/CoreStatements/HodgeTatePoints/ZarTop.lean}{PositivityTheory.HtPt.Narrow.zarhin}
